Soit $f$ la fonction définie sur $\R$ par

\[f(x)=x \e^{-2 x}.\]

On admet que $f$ est deux fois dérivable sur $\R$ et on note $f'$ la dérivée de la fonction $f$.

On note $C_{f}$ la courbe représentative de $f$ dans un repère orthonormé du plan.

\medskip

\emph{Pour chacune des affirmations suivantes, préciser si elle est vraie ou fausse, puis justifier la réponse donnée.\\Toute réponse non argumentée ne sera pas prise en compte.}

\bigskip

	\textbf{Affirmation 1.} Pour tout réel $x$, on a $f'(x)=(-2 x+1) \e^{-2 x}$.

	\bigskip

	\textbf{Affirmation 2.} La fonction $f$ est une solution sur $\R$ de l'équation différentielle :
	\[y'+2 y=\e^{-2 x}.\]

	\bigskip

	\textbf{Affirmation 3.} La fonction $f$ est convexe sur $]-\infty ~;~ 1]$.

	\bigskip

	\textbf{Affirmation 4.} L'équation $f(x)=-1$ admet une unique solution sur $\R$.

	\bigskip

	\textbf{Affirmation 5.} L'aire du domaine délimité par la courbe $C_{f}$, l'axe des abscisses et les droites d'équation $x=0$ et $x=1$ est égale à $\dfrac{1}{4}-\dfrac{3 \e^{-2}}{4}$.

\bigskip
